% Options for packages loaded elsewhere
\PassOptionsToPackage{unicode}{hyperref}
\PassOptionsToPackage{hyphens}{url}
\documentclass[
  ignorenonframetext,
  aspectratio=169,
]{beamer}
\newif\ifbibliography
\usepackage{pgfpages}
\setbeamertemplate{caption}[numbered]
\setbeamertemplate{caption label separator}{: }
\setbeamercolor{caption name}{fg=normal text.fg}
\beamertemplatenavigationsymbolsempty
% remove section numbering
\setbeamertemplate{part page}{
  \centering
  \begin{beamercolorbox}[sep=16pt,center]{part title}
    \usebeamerfont{part title}\insertpart\par
  \end{beamercolorbox}
}
\setbeamertemplate{section page}{
  \centering
  \begin{beamercolorbox}[sep=12pt,center]{section title}
    \usebeamerfont{section title}\insertsection\par
  \end{beamercolorbox}
}
\setbeamertemplate{subsection page}{
  \centering
  \begin{beamercolorbox}[sep=8pt,center]{subsection title}
    \usebeamerfont{subsection title}\insertsubsection\par
  \end{beamercolorbox}
}
% Prevent slide breaks in the middle of a paragraph
\widowpenalties 1 10000
\raggedbottom
\AtBeginPart{
  \frame{\partpage}
}
\AtBeginSection{
  \ifbibliography
  \else
    \frame{\sectionpage}
  \fi
}
\AtBeginSubsection{
  \frame{\subsectionpage}
}
\usepackage{iftex}
\ifPDFTeX
  \usepackage[T1]{fontenc}
  \usepackage[utf8]{inputenc}
  \usepackage{textcomp} % provide euro and other symbols
\else % if luatex or xetex
  \usepackage{unicode-math} % this also loads fontspec
  \defaultfontfeatures{Scale=MatchLowercase}
  \defaultfontfeatures[\rmfamily]{Ligatures=TeX,Scale=1}
\fi
\usepackage{lmodern}
\ifPDFTeX\else
  % xetex/luatex font selection
\fi
% Use upquote if available, for straight quotes in verbatim environments
\IfFileExists{upquote.sty}{\usepackage{upquote}}{}
\IfFileExists{microtype.sty}{% use microtype if available
  \usepackage[]{microtype}
  \UseMicrotypeSet[protrusion]{basicmath} % disable protrusion for tt fonts
}{}
\makeatletter
\@ifundefined{KOMAClassName}{% if non-KOMA class
  \IfFileExists{parskip.sty}{%
    \usepackage{parskip}
  }{% else
    \setlength{\parindent}{0pt}
    \setlength{\parskip}{6pt plus 2pt minus 1pt}}
}{% if KOMA class
  \KOMAoptions{parskip=half}}
\makeatother
\usepackage{longtable,booktabs,array}
\usepackage{caption}
\captionsetup[table]{skip=6pt}
\newcounter{none} % for unnumbered tables
\usepackage{calc} % for calculating minipage widths
\usepackage{caption}
% Make caption package work with longtable
\makeatletter
\def\fnum@table{\tablename~\thetable}
\makeatother
\setlength{\emergencystretch}{3em} % prevent overfull lines
\providecommand{\tightlist}{%
  \setlength{\itemsep}{0pt}\setlength{\parskip}{0pt}}
% Auto-generated by infrastructure/rendering/_pdf_latex_helpers.py::extract_math_font_preamble.
% Minimal header passed to Pandoc via -H so Beamer slides pick up the same math font
% as the combined-PDF path. Adjust the manuscript's preamble.md to change the font globally.
\usepackage{unicode-math}
\setmathfont{latinmodern-math.otf}

% Manuscript macro/package fallbacks (newcommand -> providecommand).
% Core mathematics
\usepackage{amsmath}
\usepackage{amssymb}
% Document layout
% Tables
\usepackage{booktabs}
% Caption styling (booktabs already loaded above). Small caption font with a
% bold label ("Figure 1", "Table 2") gives the math/figure-heavy layout a
% consistent, journal-like caption rhythm without fighting pandoc-crossref —
% pandoc emits standard \caption{}, and caption/captionsetup only restyle it.
% ── Theorem environments ─────────────────────────────────────────────────
% amsthm is loaded above. The FedGVI<->Friston bridge is stated as a small
% number of theorems/definitions/lemmas/propositions/corollaries; sharing one
% counter (definition/lemma/proposition/corollary all step the `theorem`
% counter) gives a single monotone numbering 1, 2, 3, ... across all kinds, so
% authors can reference them as Theorem 1, Definition 2, Corollary 3 without a
% per-kind counter drifting out of order. Plain style for theorem-like results,
% definition style (upright body) for definitions.
% (algorithm2e intentionally not loaded: this manuscript states results as
% numbered amsthm Theorems/Definitions, not algorithm2e pseudocode.)
% Code listings
\usepackage{listings}
% Typography and formatting
% (siunitx intentionally not loaded: no \SI/\num units appear in this manuscript.)
% Cross-references and citations
% (cleveref and titlesec intentionally not loaded: unavailable in this TeX tree
% and not required. Cross-references resolve through pandoc-crossref's own
% \ref-style names — no \cref appears — and section heading styling falls back
% to the pandoc/LaTeX defaults.)
% ── TeX Live Basic-compatible mono font for code listings ────────────
% Latin Modern Mono is bundled with TeX Live Basic and keeps the manuscript
% renderable on minimal TeX installations. Mathematical Unicode coverage is
% provided separately by Latin Modern Math below, so the code font does not
% need a private JuliaMono installation.
%
% Use the TeX font filename rather than the system-family name: the minimal
% TeX Live fontconfig database may not register the OpenType family even though
% kpsewhich can resolve the bundled files.
% Math font for unicode-math: Latin Modern Math (TeX Live) has full BMP
% coverage including U+2223 (\mid), U+226A/226B (\ll/\gg), and the Greek/
% blackboard letters used in equations. Without an explicit \setmathfont,
% unicode-math falls back to lmroman text font which lacks several glyphs
% and emits "Missing character" warnings on every \mid in math mode.

% Slide layout defaults for warning-clean scientific decks.
\usepackage{etoolbox}
\IfFileExists{xurl.sty}{\usepackage{xurl}}{}
\IfFileExists{seqsplit.sty}{\usepackage{seqsplit}}{\newcommand{\seqsplit}[1]{#1}}
\protected\def\breakseq#1{\seqsplit{#1}}
\protected\def\breaktt#1{\begingroup\ttfamily\seqsplit{#1}\endgroup}
\setlength{\emergencystretch}{6em}
\tolerance=5000
\hbadness=10000
\hfuzz=1pt
\setlength{\tabcolsep}{2pt}
\AtBeginEnvironment{longtable}{\tiny\renewcommand{\arraystretch}{0.86}\setlength{\tabcolsep}{1pt}}
\AtBeginEnvironment{tabular}{\tiny\renewcommand{\arraystretch}{0.86}\setlength{\tabcolsep}{1pt}}
\AtBeginEnvironment{equation}{\tiny}
\AtBeginEnvironment{equation*}{\tiny}
\AtBeginEnvironment{align}{\tiny}
\AtBeginEnvironment{align*}{\tiny}
\AtBeginEnvironment{itemize}{\footnotesize}
\AtBeginEnvironment{enumerate}{\footnotesize}
\AtBeginEnvironment{description}{\footnotesize}
\setbeamerfont{caption}{size=\tiny}
\setbeamerfont{caption name}{size=\tiny}
\setbeamerfont{normal text}{size=\small}
\setbeamerfont{frametitle}{size=\small}
\setbeamerfont{section title}{size=\footnotesize}
\setbeamerfont{subsection title}{size=\footnotesize}
\setbeamertemplate{section page}{%
  \centering
  \begin{beamercolorbox}[sep=12pt,center,wd=\paperwidth]{section title}
    \parbox{0.86\paperwidth}{\centering\usebeamerfont{section title}\insertsection\par}
  \end{beamercolorbox}
}
\setbeamertemplate{subsection page}{%
  \centering
  \begin{beamercolorbox}[sep=8pt,center,wd=\paperwidth]{subsection title}
    \parbox{0.86\paperwidth}{\centering\usebeamerfont{subsection title}\insertsubsection\par}
  \end{beamercolorbox}
}
\setlength{\abovecaptionskip}{2pt}
\setlength{\belowcaptionskip}{0pt}

% Natbib and cross-reference fallbacks — slides don't load natbib
% or cleveref, but combined-PDF manuscript prose may emit these
% commands. The fallback renders citations as a bracketed key list
% and cross-references as detokenized labels so slides don't fail on
% undefined control sequences or raw underscores. \providecommand is
% a no-op if packages load later.
\providecommand{\citep}[1]{[#1]}
\providecommand{\citet}[1]{#1}
\providecommand{\citealp}[1]{#1}
\providecommand{\citeauthor}[1]{#1}
\providecommand{\citeyear}[1]{#1}
\providecommand{\cref}[1]{\texttt{\detokenize{#1}}}
\providecommand{\Cref}[1]{\texttt{\detokenize{#1}}}
\providecommand{\crefrange}[2]{\texttt{\detokenize{#1}}--\texttt{\detokenize{#2}}}
\providecommand{\Crefrange}[2]{\texttt{\detokenize{#1}}--\texttt{\detokenize{#2}}}
\providecommand{\paragraph}[1]{\textbf{#1}\ }

% Opt-in accessible presentation profile.
\setbeamerfont{normal text}{size*={20pt}{24pt}}
\setbeamerfont{frametitle}{size*={28pt}{32pt}}
\setbeamerfont{section title}{size*={28pt}{32pt}}
\setbeamerfont{subsection title}{size*={28pt}{32pt}}
\setbeamerfont{caption}{size*={16pt}{19pt}}
\setbeamerfont{caption name}{size*={16pt}{19pt}}
\setbeamerfont{itemize/enumerate body}{size*={20pt}{24pt}}
\setbeamerfont{itemize/enumerate subbody}{size*={20pt}{24pt}}
\setbeamerfont{itemize/enumerate subsubbody}{size*={20pt}{24pt}}
\setbeamerfont{itemize item}{size*={20pt}{24pt}}
\setbeamerfont{itemize subitem}{size*={20pt}{24pt}}
\setbeamerfont{itemize subsubitem}{size*={20pt}{24pt}}
\setbeamerfont{description body}{size*={20pt}{24pt}}
\setbeamerfont{description item}{size*={20pt}{24pt}}
\setbeamerfont{quote}{size*={20pt}{24pt}}
\setbeamerfont{quotation}{size*={20pt}{24pt}}
\AtBeginEnvironment{longtable}{\fontsize{20pt}{24pt}\selectfont}
\AtBeginEnvironment{tabular}{\fontsize{20pt}{24pt}\selectfont}
\AtBeginEnvironment{equation}{\fontsize{20pt}{24pt}\selectfont}
\AtBeginEnvironment{equation*}{\fontsize{20pt}{24pt}\selectfont}
\AtBeginEnvironment{align}{\fontsize{20pt}{24pt}\selectfont}
\AtBeginEnvironment{align*}{\fontsize{20pt}{24pt}\selectfont}
\AtBeginEnvironment{itemize}{\fontsize{20pt}{24pt}\selectfont}
\AtBeginEnvironment{enumerate}{\fontsize{20pt}{24pt}\selectfont}
\AtBeginEnvironment{description}{\fontsize{20pt}{24pt}\selectfont}
\AtBeginDocument{\usebeamerfont{normal text}}
\newlength{\AFSlideTableWidth}
% The fixed footer reservation is calibrated with the semantic seven-line regular-body budget.
\setbeamertemplate{footline}{%
  \leavevmode\vbox{%
    \hbox{%
      \begin{beamercolorbox}[wd=\paperwidth,ht=16pt,dp=3pt,center]{author in head/foot}%
        \fontsize{16pt}{19pt}\selectfont Untagged PDF derivative \textbar\ \href{../web/index.html}{HTML reader}%
      \end{beamercolorbox}%
    }%
    \vskip6pt%
  }%
}
% Accessible footer reserved height: 25pt.

\newtheorem{proposition}{Proposition}
\newtheorem{hypothesis}{Hypothesis}
\newtheorem{remark}[theorem]{Remark}
\newtheorem{axiom}{Axiom}
\newtheorem{property}{Property}
\usepackage{bookmark}
\IfFileExists{xurl.sty}{\usepackage{xurl}}{} % add URL line breaks if available
\urlstyle{same}
% fallback for those not using the hyperref driver hyperxmp:
\makeatletter
\@ifundefined{xmpquote}{\newcommand{\xmpquote}[1]{#1}}{}
\makeatother
\hypersetup{
  hidelinks,
  pdfcreator={LaTeX via pandoc}}

\author{\texorpdfstring{}{}}
\date{}

\begin{document}

\section{Supplemental notation contract}\label{sec:supp-notation}

\begin{frame}{Supplemental notation contract}
This supplement is the authoritative notation contract for Active
Fedference. The methods, formalism, results, figures, report schemas,
and API documentation use these meanings even when a source paper uses a
different symbol.
\end{frame}

\begin{frame}{Supplemental notation\ldots{} (part 2)}
A symbol is not reused for a different mathematical object merely
because the objects are both probability vectors.
\end{frame}

\begin{frame}{Supplemental notation\ldots{} (part 3)}
The implementation names in the final column are the canonical names for
new code and reports; old names survive only as warned, parity-tested
compatibility adapters.
\end{frame}

\section{Probability objects and generative-model
quantities}\label{probability-objects-and-generative-model-quantities}

\begin{frame}{States, posteriors, and site factors}
\protect\phantomsection\label{states-posteriors-and-site-factors}
{\def\LTcaptype{none} % do not increment counter
% template-accessible-table-inset
\begingroup
\setlength{\AFSlideTableWidth}{\textwidth}%
\setlength{\LTleft}{\fill}%
\setlength{\LTright}{\fill}%
\begin{longtable}[]{@{}
  >{\raggedright\arraybackslash}p{(\AFSlideTableWidth - 2\tabcolsep) * \real{0.2963}}
  >{\raggedright\arraybackslash}p{(\AFSlideTableWidth - 2\tabcolsep) * \real{0.7037}}@{}}
\toprule\noalign{}
\begin{minipage}[b]{\linewidth}\raggedright
Symbol
\end{minipage} & \begin{minipage}[b]{\linewidth}\raggedright
Meaning
\end{minipage} \\
\midrule\noalign{}
\endhead
\(s\) & Hidden categorical-state index. \\
\(o\) & Observation/outcome index. \\
\(q_n(s)\) & Agent \(n\)'s updated local posterior. \\
\(q(s)\) & Server consensus posterior. \\
\bottomrule\noalign{}
\end{longtable}
\endgroup
}
\end{frame}

\begin{frame}{States, posteriors, and site\ldots{} (part 2)}
Here \(s\in\{1,\ldots,n_s\}\) and \(o\in\{1,\ldots,n_o\}\) are indices,
not distributions. Each \(q_n(s)\) is the local posterior over the
shared latent state after agent \(n\)'s update; \(q(s)\) is the
corresponding server consensus, and \(q_{-n}(s)\) removes agent \(n\)'s
site contribution. The site \(t_n(s)\) is a natural-parameter factor.
\end{frame}

\begin{frame}{States, posteriors, and site\ldots{} (part 3)}
The bridge-only quantity satisfies \(m_n(s)=\log q_n(s)+\kappa_n\) for
state-constant \(\kappa_n\); it has no project API and does not imply
that every source-protocol message is a broadcast posterior.
\end{frame}

\begin{frame}[fragile]{States, posteriors, and site\ldots{} (part 4)}
The corresponding code terms are \texttt{state}, \texttt{observation},
\texttt{local\_posteriors{[}n{]}}, \texttt{global\_posterior} or
\texttt{consensus}, \texttt{cavity(...)}, and \texttt{site\_factor}, in
table order through \(t_n(s)\). The bridge-only \(m_n(s)\) has no
project API.
\end{frame}

\begin{frame}{Priors, policies, and POMDP quantities}
\protect\phantomsection\label{priors-policies-and-pomdp-quantities}
{\def\LTcaptype{none} % do not increment counter
% template-accessible-table-inset
\begingroup
\setlength{\AFSlideTableWidth}{\textwidth}%
\setlength{\LTleft}{\fill}%
\setlength{\LTright}{\fill}%
\begin{longtable}[]{@{}
  >{\raggedright\arraybackslash}p{(\AFSlideTableWidth - 2\tabcolsep) * \real{0.5789}}
  >{\raggedright\arraybackslash}p{(\AFSlideTableWidth - 2\tabcolsep) * \real{0.4211}}@{}}
\toprule\noalign{}
\begin{minipage}[b]{\linewidth}\raggedright
Symbol
\end{minipage} & \begin{minipage}[b]{\linewidth}\raggedright
Meaning
\end{minipage} \\
\midrule\noalign{}
\endhead
\(\pi_0(s)\) & Hidden-state prior, not a policy. \\
\bottomrule\noalign{}
\end{longtable}
\endgroup
}
\end{frame}

\begin{frame}{Priors, policies, and\ldots{} (part 2)}
The state index \(s\), posterior \(q(s)\), prior \(\pi_0(s)\), and
policy \(\boldsymbol{\pi}\) must not be conflated. In particular, the
policy symbol is bold when needed, never used for a prior, and the prior
is never called a policy. The uppercase POMDP tensors \(A,B,C,D_0\) are
model objects; they are not posterior factors.
\end{frame}

\begin{frame}{Priors, policies, and\ldots{} (part 3)}
Each state-indexed column of \(A\) is a pmf. The tensor \(B\) is indexed
by next state, current state, and control. The preferred-outcome pmf is
\(p_C(o)=\operatorname{softmax}(C)[o]\), and
\(q(o\mid\boldsymbol{\pi})\) is the predicted-outcome distribution used
in EFE calculations.
\end{frame}

\begin{frame}[fragile]{Priors, policies, and\ldots{} (part 4)}
The canonical code terms, in table order, are \texttt{prior} or
\texttt{log\_prior}, \texttt{policy}, \texttt{likelihood},
\texttt{transition}, \texttt{log\_preferences}, \texttt{initial\_prior},
and \texttt{predicted\_outcomes}.
\end{frame}

\begin{frame}{Priors, policies, and\ldots{} (part 5)}
For the qualified relation to Friston et al.'s Eq. 7 (Friston et al.
2024), the shared support is finite, every \(q_n(s)\) is positive on it,
the Eq. 7 softmax input is represented by the bridge-only \(m_n(s)\)
above, and the declared weights \(w_n\) are fixed rather than functions
of the emerging consensus.
\end{frame}

\begin{frame}{Priors, policies, and\ldots{} (part 6)}
Under exactly those assumptions, additive \(\kappa_n\) constants cancel
under softmax and eq.~9 is the categorical
posterior-log-potential specialization of the source message-combination
term. It neither reconstructs source message construction,
cavity/exclusion policy, scheduling,
\end{frame}

\begin{frame}{Priors, policies, and\ldots{} (part 7)}
generative factors, nor the complete source protocol.
\end{frame}

\section{Divergences, losses, and scalar
controls}\label{divergences-losses-and-scalar-controls}

\begin{frame}{Generalized-Bayes and aggregation terms}
\protect\phantomsection\label{generalized-bayes-and-aggregation-terms}
{\def\LTcaptype{none} % do not increment counter
% template-accessible-table-inset
\begingroup
\setlength{\AFSlideTableWidth}{\textwidth}%
\setlength{\LTleft}{\fill}%
\setlength{\LTright}{\fill}%
\begin{longtable}[]{@{}
  >{\raggedright\arraybackslash}p{(\AFSlideTableWidth - 2\tabcolsep) * \real{0.5000}}
  >{\raggedright\arraybackslash}p{(\AFSlideTableWidth - 2\tabcolsep) * \real{0.5000}}@{}}
\toprule\noalign{}
\begin{minipage}[b]{\linewidth}\raggedright
Symbol
\end{minipage} & \begin{minipage}[b]{\linewidth}\raggedright
Meaning
\end{minipage} \\
\midrule\noalign{}
\endhead
\(\mathcal D(q\Vert p)\) & Regularizing divergence. \\
\bottomrule\noalign{}
\end{longtable}
\endgroup
}
\end{frame}

\begin{frame}[fragile]{Generalized-Bayes and\ldots{} (part 2)}
The regularizing divergence acts between distributions and may be KL,
reverse KL, or \(\alpha\)-Rényi. The positive learning rate/temperature
\(\tau\) multiplies accumulated loss and accepts \texttt{learning\_rate}
only as a warned compatibility alias.
\end{frame}

\begin{frame}{Generalized-Bayes and\ldots{} (part 3)}
The symbols \(w_n\), \(a_n\), and \(\widetilde a_n\) are deliberately
distinct. The first is supplied before aggregation, the second is the
variational raw server output satisfying \(0\le a_n\le w_n\), and the
third is only its normalized influence representation returned for
interpretation and plotting.
\end{frame}

\begin{frame}[fragile]{Generalized-Bayes and\ldots{} (part 4)}
Their canonical code terms, in table order, are \texttt{divergence},
\texttt{loss\_by\_state}, \texttt{tau}, \texttt{base\_weights{[}n{]}},
\texttt{raw\_effective\_weights{[}n{]}}, and
\texttt{normalized\_effective\_weights{[}n{]}}.
\end{frame}

\begin{frame}[fragile]{Generalized-Bayes and\ldots{} (part 5)}
The server heuristic's reweighting is not a FedGVI client loss and does
not inherit a client-side robustness theorem. \texttt{robust\_aggregate}
is a server heuristic with the tested \(c=0\) recovery identity.
\end{frame}

\begin{frame}[fragile]{Generalized-Bayes and\ldots{} (part 6)}
\texttt{variational\_aggregate} owns the explicit finite-simplex
objective and the raw-weight bound; that bound is not an estimator-level
B-robustness theorem.
\end{frame}

\begin{frame}{Robustness, divergences, and loss controls}
\protect\phantomsection\label{robustness-divergences-and-loss-controls}
{\def\LTcaptype{none} % do not increment counter
% template-accessible-table-inset
\begingroup
\setlength{\AFSlideTableWidth}{\textwidth}%
\setlength{\LTleft}{\fill}%
\setlength{\LTright}{\fill}%
\begin{longtable}[]{@{}
  >{\raggedright\arraybackslash}p{(\AFSlideTableWidth - 2\tabcolsep) * \real{0.3929}}
  >{\raggedright\arraybackslash}p{(\AFSlideTableWidth - 2\tabcolsep) * \real{0.6071}}@{}}
\toprule\noalign{}
\begin{minipage}[b]{\linewidth}\raggedright
Symbol
\end{minipage} & \begin{minipage}[b]{\linewidth}\raggedright
Meaning
\end{minipage} \\
\midrule\noalign{}
\endhead
\(c\ge0\) & Server divergence-reweighting coefficient. \\
\bottomrule\noalign{}
\end{longtable}
\endgroup
}
\end{frame}

\begin{frame}{Robustness, divergences\ldots{} (part 2)}
The coefficient \(\lambda\) controls the variational objective and its
coordinate updates; \(\lambda\downarrow0\) is a separate deterministic
tied-argmax endpoint. The parameter \(q_{\text{loss}}\) belongs to
\(L_{q_{\text{loss}}}\);
\end{frame}

\begin{frame}{Robustness, divergences\ldots{} (part 3)}
its \(q_{\text{loss}}\downarrow0\) NLL limit is handled separately, and
the subscript prevents collision with posterior \(q(s)\). Contamination
rate \(\rho\) is always interpreted within the declared attack
mechanism.
\end{frame}

\begin{frame}[fragile]{Robustness, divergences\ldots{} (part 4)}
The associated code terms are \texttt{robustness},
\texttt{entropy\_weight}, \texttt{alpha}, \texttt{beta},
\texttt{q\_loss}, and \texttt{contamination\_rate} or \texttt{rate}, in
table order.
\end{frame}

\begin{frame}{Robustness, divergences\ldots{} (part 5)}
For the objective-backed server rule, the complete scalar-control
contract is defined for \(c>0\) and \(\lambda>0\):
\end{frame}

\begin{frame}{Robustness, divergences\ldots{} (part 6)}
\begin{equation}\tag{34}\protect\phantomsection\label{eq:notation-variational-objective}{
\begin{aligned}
F_\lambda(q,a)
 &= \sum_n a_n\,\mathrm{CE}(q,q_n)
    -\lambda H(q)\\
 &\qquad +\frac{1}{c}\,\mathrm{KL}_{\rm gen}(a\Vert w).
\end{aligned}
}\end{equation}
\end{frame}

\begin{frame}{Robustness, divergences\ldots{} (part 7)}
Its two coordinate updates are shown separately so each display remains
legible at the presentation typography floor:
\end{frame}

\begin{frame}{Robustness, divergences\ldots{} (part 8)}
\begin{equation}\tag{35}\protect\phantomsection\label{eq:notation-variational-updates}{
\begin{aligned}
q &\propto \exp\!\left(
\frac{1}{\lambda}\sum_n a_n\log q_n\right),\\
a_n &= w_n\exp\!\big[-c\,\mathrm{CE}(q,q_n)\big].
\end{aligned}
}\end{equation}
\end{frame}

\begin{frame}[fragile]{Robustness, divergences\ldots{} (part 9)}
The coordinate updates eq.~35 use
\(\lambda=1.0\) by default; the \texttt{entropy\_weight} argument
exposes the stated tempered family. The \(c=0\) branch is handled as a
recovery limit outside the \(c>0\) objective; at \(\lambda=1.0\) it is
the exact project log-linear pool.
\end{frame}

\begin{frame}{Robustness, divergences\ldots{} (part 10)}
The \(\lambda\downarrow0\) endpoint is separately implemented as a
deterministic tied-argmax rule and is not obtained by substituting
\(\lambda=0\) into the displayed objective or update.
\end{frame}

\begin{frame}{Cavity and factor algebra}
\protect\phantomsection\label{cavity-and-factor-algebra}
For a positive-support global posterior and site term, the cavity
operation is defined in log space and then normalized:
\end{frame}

\begin{frame}{Cavity and factor algebra (part 2)}
\begin{equation}\tag{36}\protect\phantomsection\label{eq:notation-cavity}{
\begin{aligned}
q_{-n}(s)
&= \frac{q(s)/t_n(s)}{\sum_{s'}q(s')/t_n(s')}\\
&= \operatorname{softmax}\!\left(\log q(s)-\log t_n(s)\right).
\end{aligned}
}\end{equation}
\end{frame}

\begin{frame}{Cavity and factor algebra (part 3)}
The corresponding factor replacement is
\end{frame}

\begin{frame}{Cavity and factor algebra (part 4)}
\begin{equation}\tag{37}\protect\phantomsection\label{eq:notation-factor-replacement}{
\begin{aligned}
\ell_n(s)
&= \log t_n^{\mathrm{old}}(s)\\
&\quad +\log q^{\mathrm{new}}(s)-\log q^{\mathrm{old}}(s),\\
t_n^{\mathrm{new}}(s)
&\leftarrow \operatorname{softmax}(\ell_n)(s).
\end{aligned}
}\end{equation}
\end{frame}

\begin{frame}{Cavity and factor algebra (part 5)}
Here \(\ell_n\) is the transient, unnormalized log-site vector; the
softmax is exactly the exponential normalization of that vector.
\end{frame}

\begin{frame}[fragile]{Cavity and factor algebra (part 6)}
The code-level \texttt{cavity} adapter takes \texttt{global\_posterior}
and \texttt{site\_factor}. The \texttt{update\_factor} adapter takes
\texttt{old\_site\_factor}, \texttt{old\_global\_posterior}, and
\texttt{new\_global\_posterior}. The old keywords \texttt{posterior},
\texttt{factor}, \texttt{old\_factor}, \texttt{old\_posterior},
\end{frame}

\begin{frame}[fragile]{Cavity and factor algebra (part 7)}
and \texttt{new\_posterior} are accepted only with a
\texttt{DeprecationWarning}; mixed canonical/old calls fail closed.
\end{frame}

\begin{frame}{Cavity and factor algebra (part 8)}
Recombination is tested by normalizing \(q_{-n}(s)t_n(s)\) and checking
recovery of \(q(s)\) to floating-point tolerance. A transported site
factor is represented as a pmf, so its arbitrary positive
natural-parameter scale is fixed by the explicit normalization above.
\end{frame}

\begin{frame}{Statistical notation and nesting}
\protect\phantomsection\label{statistical-notation-and-nesting}
{\def\LTcaptype{none} % do not increment counter
% template-accessible-table-inset
\begingroup
\setlength{\AFSlideTableWidth}{\textwidth}%
\setlength{\LTleft}{\fill}%
\setlength{\LTright}{\fill}%
\begin{longtable}[]{@{}
  >{\raggedright\arraybackslash}p{(\AFSlideTableWidth - 2\tabcolsep) * \real{0.3000}}
  >{\raggedright\arraybackslash}p{(\AFSlideTableWidth - 2\tabcolsep) * \real{0.7000}}@{}}
\toprule\noalign{}
\begin{minipage}[b]{\linewidth}\raggedright
Symbol
\end{minipage} & \begin{minipage}[b]{\linewidth}\raggedright
Contractual meaning
\end{minipage} \\
\midrule\noalign{}
\endhead
\(n_{\rm seed}\) & Number of independently seeded worlds/replicates in a
declared cell; the inferential unit for seed-level summaries. \\
\bottomrule\noalign{}
\end{longtable}
\endgroup
}
\end{frame}

\begin{frame}{Statistical notation and\ldots{} (part 2)}
For the robustness sweep, the primary result is \(r_{\rm rb}\) and the
matched mean difference \(\overline{\Delta}\) with its bootstrap
interval. The d-equivalent is retained only as a monotone secondary
display and planning input. When \(|r_{\rm rb}|=1\), the transform
diverges;
\end{frame}

\begin{frame}{Statistical notation and\ldots{} (part 3)}
reports use a finite sentinel and captions disclose saturation rather
than presenting a million-scale number as a scientifically interpretable
effect.
\end{frame}

\begin{frame}{Statistical notation and\ldots{} (part 4)}
Power, prospective sample size, MCSE, and MDE are observed-effect
planning/precision diagnostics, not evidence that a confirmatory effect
exists.
\end{frame}

\begin{frame}{Statistical notation and\ldots{} (part 5)}
The predeclared headline display rule is the largest positive
\(r_{\rm rb}\) among robust methods, with the declared method order as a
deterministic tie break. A report must also expose the complete
tied-method set, the tie-break, the method with the largest mean
\(\overline{\Delta}\),
\end{frame}

\begin{frame}{Statistical notation and\ldots{} (part 6)}
and the method with the largest mean at the worst rate.

These are distinct summaries; none is a unique scientific winner when
the evidence is tied or conditional.
\end{frame}

\begin{frame}{Statistical notation and\ldots{} (part 7)}
For the review grid, every configured robust method remains an
inferential member and a displayed rate-profile curve. No pooled-mean
selection creates a curve, interval, or hypothesis-test member for that
surface.
\end{frame}

\begin{frame}[fragile]{Code and manuscript naming map}
\protect\phantomsection\label{code-and-manuscript-naming-map}
Every arrow names a legacy surface followed by its canonical
replacement.

\textbf{Belief inputs.} \texttt{beliefs} or \texttt{agent\_beliefs} \texorpdfstring{\ensuremath{\to}}{->}
\texttt{local\_posteriors}; \texttt{weights} \texorpdfstring{\ensuremath{\to}}{->} \texttt{base\_weights}.
Both are warned keyword/property adapters.
\end{frame}

\begin{frame}[fragile]{Code and manuscript\ldots{} (part 2)}
\textbf{Aggregation outputs.} Result property \texttt{agent\_weights} \texorpdfstring{\ensuremath{\to}}{->}
\texttt{normalized\_effective\_weights}; variational argument
\texttt{agent\_weights} \texorpdfstring{\ensuremath{\to}}{->} \texttt{raw\_effective\_weights}. Both are
warned adapters, but represent different weight scales.
\end{frame}

\begin{frame}[fragile]{Code and manuscript\ldots{} (part 3)}
\textbf{Diagnostics.} \texttt{shared\_beliefs} \texorpdfstring{\ensuremath{\to}}{->}
\texttt{shared\_posteriors}; \texttt{loss\_vec} \texorpdfstring{\ensuremath{\to}}{->}
\texttt{loss\_by\_state}. These are warned diagnostics and keyword
adapters.
\end{frame}

\begin{frame}[fragile]{Code and manuscript\ldots{} (part 4)}
\textbf{Statistics.} \texttt{cohens\_d\_from\_rank\_biserial} \texorpdfstring{\ensuremath{\to}}{->}
\texttt{d\_equivalent\_from\_rank\_biserial}; report field
\texttt{cohens\_d} \texorpdfstring{\ensuremath{\to}}{->} \texttt{d\_equivalent}. The function is a warned
adapter; the report change is a versioned migration.
\end{frame}

\begin{frame}[fragile]{Code and manuscript\ldots{} (part 5)}
These adapters never silently reinterpret an argument or result. The
\texttt{agent\_weights} result property preserves serialized wire
compatibility, while the variational objective and new reports use the
canonical weight terms. New reports are canonical and schema-versioned;
\end{frame}

\begin{frame}[fragile]{Code and manuscript\ldots{} (part 6)}
readers fail closed on unsupported versions.
\texttt{d\_equivalent\_from\_rank\_biserial} returns the declared
rank-biserial-derived d-equivalent, not raw Cohen's \(d\);
\texttt{loss\_vec} remains only a warned generalized-Bayes keyword
adapter.
\end{frame}

\begin{frame}[fragile]{Code and manuscript\ldots{} (part 7)}
The wire-level key \texttt{agent\_weights} is preserved because it is a
federation transport contract, not a claim about the scale or meaning of
the new result fields. A future wire migration requires an explicit
version and a fail-closed reader; it must not silently reinterpret the
key.
\end{frame}

\begin{frame}{Source and evidence boundaries}
\protect\phantomsection\label{source-and-evidence-boundaries}
The Friston belief-sharing equations are source equations/protocol
claims. Only under the explicit finite-shared-support,
posterior-log-potential, and fixed-weight bridge above does the
categorical log-linear pool specialize the source message-combination
term;
\end{frame}

\begin{frame}{Source and evidence\ldots{} (part 2)}
the tested \(c=0\) identity remains project-local. The generalized-Bayes
and loss limits are implementation analogues checked in the finite
categorical model.
\end{frame}

\begin{frame}{Source and evidence\ldots{} (part 3)}
The contamination, gallery, onset, conditional-world, and review-grid
quantities are conditional simulation evidence over declared cells. None
of these finite surfaces is an external-data replication, a
reconstructed source protocol, a universal attack taxonomy, a causal
intervention,
\end{frame}

\begin{frame}{Source and evidence\ldots{} (part 4)}
or a proof of a server-side robustness guarantee.
\end{frame}

\begin{frame}{Source and evidence\ldots{} (part 5)}
Open theory, calibration, protocol, continuous-state, external-data,
authenticated-federation, and clean-release work remains open in the
project TODO and claim-audit documents.
\end{frame}

\end{document}
